\section{A finite inequality for the drift}
\label{sec:inequality}

The cost analysis uses an inequality on a finite probability space.  Its
comparison matrix need not be symmetric.  The proof allows a positive
contribution at one index only by charging indices with smaller values;
the negative contributions at those indices pay the entire charge.

The statement below is independent of the rounding construction.
Section~\ref{sec:rounding} will apply it to normalized client distances
and directed edge probabilities; Equation~\ref{eq:rnd-normalization}
records that substitution. We prove the finite statement first so that
the geometric analysis can use it as a single input.

For $0\le x\le1$, define
\begin{equation}
\label{eq:analytic-coefficients}
\begin{aligned}
 A(x)&=1+2\int_0^1(1-e^{-tx})\,dt,&
 P(x)&=\int_0^1te^{-tx}\,dt,\\
 B(x)&=2P(x),& m(x)&=A(x)-1-2xP(x).
\end{aligned}
\end{equation}
We suppress the argument $x$ when it is fixed.  Differentiation under the
integral gives $A'=B$ and
$A''(x)=-2\int_0^1t^2e^{-tx}\,dt<0$.  In particular,
$0<P\le1/2$ and $m\ge0$, the latter by concavity and $A(0)=1$.
Integration by parts also gives
\begin{equation}
\label{eq:analytic-identities}
 (1+x)(A-1)=2x(1-P),\qquad
 (1+x)A+xB-1=3x,\qquad A(1)=1+\frac2e.
\end{equation}
These identities hold at $x=0$ as well, using the integral definitions.

\begin{theorem}[Finite comparison inequality]
\label{thm:finite-inequality}
Let $I$ be a finite index set with probability weights $(\pi_i)_{i\in I}$,
and let $D_i\ge0$ satisfy $\E_iD_i=1$.  Both $\E_i$ and $\E_t$ denote
expectation with these weights; in a double expectation the two indices
are sampled independently.  Fix $0\le x\le1$ and $0<y\le1$, and set
\[
 \Theta=\frac{P(y-x)^2}{y}.
\]
Suppose $q_{ti}\in[0,1]$ and $R_i\ge0$ satisfy
\begin{equation}
\label{eq:finite-comparison}
 R_i(1-Pq_{ti})\le D_i+D_t\qquad\text{whenever }q_{ti}<1.
\end{equation}
No symmetry is assumed for $q$.  For any $s_i\in[0,1]$, put
$r_i=\E_tq_{ti}$ and $C_i=1-Ps_i$.  Then
\begin{equation}
\label{eq:finite-inequality}
\begin{aligned}
 \E_i\Big[&C_i(1-r_i)R_i
 +r_i\{mD_i+(2-s_i)Py(D_i+1)\}-C_i(D_i+1)\\
 &\hspace{12mm}
 -\E_t(1-q_{ti})\bigl(D_t-(1+m+2Py)D_i\bigr)_+\Big]
 \le2\Theta.
\end{aligned}
\end{equation}
Here $u_+=\max(u,0)$.
\end{theorem}

The proof will bound each positive row by charges sent to indices with
smaller $D$ values, then show that those receiving rows pay all incoming
charges. The scalar estimates below locate the positive rows, bound the
largest value they can charge, and certify the required negative
contributions. Every terminating decimal denotes an exact rational number.

\begin{lemma}[Scalar bounds]
\label{lem:scalar-bounds}
Fix $x\in[0,1]$, put $X=1-x$, and take
$0\le w\le X$ and $s\in\{0,1\}$.  Use the functions in
Equation~\eqref{eq:analytic-coefficients}, and define
\begin{equation}
\label{eq:scalar-parameters}
\begin{aligned}
 z&=Pw^2,& b&=(2-s)P(x+w),& a&=A+2Pw,& C&=1-Ps,\\
 d_0&=.52,& d_1&=.4,& H_0&=.394-.42X^2+.04X^3,& H_1&=.25,\\
 J_s&=CH_s,&
 k_s&=(1-s)\frac{1-\sqrt{1-P}}{1+\sqrt{1-P}}.
\end{aligned}
\end{equation}
For $q,r\in[0,1]$, set
\begin{equation}
\label{eq:scalar-functions}
\begin{aligned}
 \psi(q)&=\frac{1-q}{1-Pq},&
 \kappa_s(q)&=C\psi(q)-(1-q),\\
 K(r)&=C-k_s-br+1.6z,&
 L(r)&=(C-m-b)r+.4z,\\
 M(r)&=(1-r)a+CPr,& G(r)&=M(r)-L(r).
\end{aligned}
\end{equation}
Then the following five bounds hold on the indicated domains:
\begin{align}
 m+b-\frac{.4z}{r}
 &\le\frac{C(1-P)}{1-Pr}
 &&(0<r\le1),
 \label{eq:scalar-first}\\
 L(r)\ge0,\qquad K(r)>0,\qquad K(r)
 &\ge d_sG(r)
 &&(0\le r\le1),
 \label{eq:scalar-linear}\\
 \bigl((J_s-b)r+1.6z\bigr)M(r)
 +(C-k_s-J_sr)L(r)
 &\ge0
 &&(0\le r\le1),
 \label{eq:scalar-charge}\\
 C-k_s+1.6z-bq+.52\bigl(k_s-\kappa_s(q)\bigr)
 &\ge H_0
 &&(0\le q\le1),
 \label{eq:scalar-type-zero}\\
 C-k_s+1.6z-b&\ge.25.
 \label{eq:scalar-type-one}
\end{align}
\end{lemma}

\begin{proof}
The scalar bounds are proved in Appendix~\ref{app:certificates} by
analytic reductions and exact polynomial certificates.
\end{proof}

\begin{proof}[Proof of Theorem~\ref{thm:finite-inequality}]
We may discard indices of probability zero.  The proof treats the
quantities with a fixed second index $i$ as a row, with $t$ indexing its
columns.

\paragraph{Reduction of the parameters.}
For fixed $x,D,q,R,s$, the left side of
Equation~\eqref{eq:finite-inequality} is nondecreasing in $y$.
The coefficient of its linear $y$ term is nonnegative, and increasing
$y$ decreases the positive part that is subtracted.  Replace $y$ there
by $\max(x,y)=x+w$, where $0\le w\le X=1-x$, and put $z=Pw^2$.
If $y<x$, then $z=0\le\Theta$; otherwise
$z=P(y-x)^2\le\Theta$, since $y\le1$.  It is therefore sufficient to
bound the modified left side by $2z$.

This modified left side is affine separately in every $s_i$, and its
hypotheses do not constrain $s_i$.  Its maximum over $[0,1]^I$ is
attained with every $s_i\in\{0,1\}$.  Fix such a choice.  For each row
use the notation of Lemma~\ref{lem:scalar-bounds}, with $s=s_i$, and
temporarily suppress the row subscript.  Write $D=D_i$, $R=R_i$,
$q=q_{ti}$, $Z=D_t$, and $r=\E_tq_{ti}$, and put
\[
 U=mD+b(D+1).
\]
Notice that $a=1+m+2P(x+w)$.  Define the row contribution
\begin{equation}
\label{eq:row-contribution}
 F=C(1-r)R+Ur-C(D+1)
       -\E_t(1-q)(Z-aD)_+-(.4D+1.6)z.
\end{equation}
Because $\E_iD_i=1$, proving $\E_iF_i\le0$ proves the desired
bound $2z$.

\paragraph{A centered majorant for each row.}
For $s=0$, direct maximization gives
\[
 \kappa_0(q)=\frac{Pq(1-q)}{1-Pq}
 \le\frac{1-\sqrt{1-P}}{1+\sqrt{1-P}}=k_0.
\]
For $s=1$, one has
$\kappa_1(q)=-P(1-q)^2/(1-Pq)\le0=k_1$.
Also
\begin{equation}
\label{eq:psi-bound}
 \psi(q)=1-q+Pq\psi(q)\le1-q+Pq,
\end{equation}
since $\psi(q)\le1$.

Equation~\eqref{eq:finite-comparison} implies
$(1-q)R\le\psi(q)(D+Z)$ when $q<1$.  This last inequality also holds
when $q=1$, as both sides are zero.  Using
$(Z-aD)_+\ge Z-aD$ in Equation~\eqref{eq:row-contribution} gives
\begin{align*}
 F\le\E_t\big[&Uq+D\{C\psi(q)+a(1-q)-C\}
              -(.4D+1.6)z-C+\kappa_s(q)Z\big].
\end{align*}
The centering identity
\[
 \E_t[\kappa_s(q)Z]
 =k_s-\E_t[(k_s-\kappa_s(q))Z]
\]
uses $\E_tZ=1$.  Thus, defining the explicit function
\begin{equation}
\label{eq:centered-integrand}
\begin{aligned}
 I_s(D;q,Z)={}&[mD+b(D+1)]q
 +D\{C\psi(q)+a(1-q)-C\}\\
 &-(.4D+1.6)z-C+k_s-(k_s-\kappa_s(q))Z,
\end{aligned}
\end{equation}
we have
\begin{equation}
\label{eq:centered-row-bound}
 F_i\le\E_t I_{s_i}(D_i;q_{ti},D_t).
\end{equation}
The replacement of $k_sZ$ by $k_s$ is made only under expectation.
Subsequent pointwise bounds will concern the function in
Equation~\eqref{eq:centered-integrand}.

Since $k_s-\kappa_s(q)\ge0$ and $Z\ge0$, dropping the last term and
using Equation~\eqref{eq:psi-bound} yields
\begin{equation}
\label{eq:pointwise-linear-majorant}
 I_s(D;q,Z)\le DG(q)-K(q).
\end{equation}
Both $G$ and $K$ are affine in their argument, so averaging gives
\begin{equation}
\label{eq:row-linear-majorant}
 F\le DG(r)-K(r).
\end{equation}

\paragraph{Charges issued by positive rows.}
If $r=0$, then $q=0$ at every column and averaging
Equation~\eqref{eq:finite-comparison} gives $R\le D+1$.
Equation~\eqref{eq:row-contribution} then gives $F\le0$.
Consequently a row with $F>0$ has $r>0$.  By
Equations~\eqref{eq:scalar-linear} and
\eqref{eq:row-linear-majorant}, it also has
\begin{equation}
\label{eq:positive-row-location}
 G(r)>0,\qquad D>\frac{K(r)}{G(r)}\ge d_s.
\end{equation}
In this row put
\[
 U'=\frac{Ur-(.4D+1.6)z}{Cr}.
\]
The identity
\[
 U'=\frac{D+1}{C}\left(m+b-\frac{.4z}{r}\right)
       -\frac mC-\frac{1.2z}{Cr}
\]
and Equation~\eqref{eq:scalar-first} imply
\begin{equation}
\label{eq:positive-row-uprime}
 U'\le\frac{(D+1)(1-P)}{1-Pr}.
\end{equation}
On the other hand, discarding the nonpositive hinge term in
Equation~\eqref{eq:row-contribution} gives
$F/C\le(1-r)R+rU'-(D+1)$.
If $R\le(D+1)/(1-Pr)$, Equation~\eqref{eq:positive-row-uprime}
makes this upper bound nonpositive.  Therefore
\begin{equation}
\label{eq:positive-row-radius}
 R>\frac{D+1}{1-Pr},\qquad U'\le R(1-P).
\end{equation}
In particular $r=1$ is impossible for a positive row, since
Equation~\eqref{eq:positive-row-uprime} would then give $U'\le D+1$.

For every column with $q<1$, Equations~\eqref{eq:finite-comparison}
and \eqref{eq:positive-row-radius} now give
\[
 (1-q)R+qU'\le R(1-Pq)\le D+Z.
\]
For $q=1$, the excess of the left side over $D+Z$ is exactly
$U'-D-Z$.  Using $\E_tZ=1$ and again discarding the nonpositive
hinge term, we obtain
\begin{equation}
\label{eq:positive-row-excess}
 F\le C\E_t\big[\ind_{\{q=1\}}(U'-D-Z)_+\big]
   \le C\E_t(U'-D-Z)_+.
\end{equation}
This treats the unconstrained columns with $q=1$ explicitly.

It remains to bound the largest value to which this row can charge.
An exact rearrangement gives
\[
 U'-D=\frac{br-1.6z-L(r)D}{Cr}.
\]
Set $T=(J_s-b)r+1.6z$.  The expression in
Equation~\eqref{eq:scalar-charge} is
\begin{equation}
\label{eq:charge-scalar-rearrangement}
 TM(r)+(C-k_s-J_sr)L(r)=TG(r)+K(r)L(r)\ge0.
\end{equation}
Using $L(r)\ge0$ and Equation~\eqref{eq:positive-row-location},
we conclude that
\[
 br-1.6z-L(r)D
 \le br-1.6z-\frac{L(r)K(r)}{G(r)}\le J_sr.
\]
Thus $U'-D\le H_s$.  Since $C\le1$,
Equation~\eqref{eq:positive-row-excess} proves
\begin{equation}
\label{eq:outgoing-charge}
 F_i\le\E_t(H_{s_i}-D_t)_+\qquad(F_i>0).
\end{equation}
The function $H_0$ decreases from $.394$ to $.014$ on $[0,1]$,
and $H_1=.25$.  Consequently every nonzero charge in
Equation~\eqref{eq:outgoing-charge} goes to an index with $D_t<.394$.
This separates every possible receiving index from all positive rows,
which have $D_i>d_{s_i}\ge.4$. It remains to prove that the negative
contribution of each receiving row covers the total charge it receives.

\paragraph{Negative rows cover the incoming charges.}
Put $c_*=.394$ and define
\[
 \mathcal P=\{i:F_i>0\},\qquad
 \mathcal S=\{i:D_i\le c_*\},\qquad
 \Gamma_{it}=\ind_{\{i\in\mathcal P\}}(H_{s_i}-D_t)_+.
\]
Equation~\eqref{eq:positive-row-location} shows that
$\mathcal P\cap\mathcal S=\varnothing$, because $d_s\ge.4>c_*$.
We claim that every $i\in\mathcal S$ satisfies
\begin{equation}
\label{eq:incoming-charge}
 F_i+\E_t\Gamma_{ti}\le0.
\end{equation}

Fix such an $i$ and a column $t$.  Hold $q=q_{ti}$, $Z=D_t$,
the two types $s_i,s_t$, and $\chi=\ind_{\{F_t>0\}}$ fixed.
As a function of a variable $D\in[0,c_*]$,
\[
 I_{s_i}(D;q,Z)+\chi(H_{s_t}-D)_+
\]
is convex: the first term is affine and the second is a positive part
multiplied by a nonnegative constant.  It suffices to bound this
expression at $D=0,c_*$.  This check concerns the explicit majorant;
it does not require the comparison hypothesis to remain valid when
the variable $D$ changes.

At $D=c_*$, the incoming charge is zero.  By
Equation~\eqref{eq:pointwise-linear-majorant}, the remaining term is
at most $c_*G(q)-K(q)\le0$.  Indeed, if $G(q)<0$, use $K(q)>0$;
otherwise use $K(q)\ge d_{s_i}G(q)$ and $c_*<d_{s_i}$.

At $D=0$, Equation~\eqref{eq:centered-integrand} reads
\[
 I_{s_i}(0;q,Z)
 =bq-1.6z-C+k_s-(k_s-\kappa_s(q))Z,
\]
where the unindexed scalar parameters belong to row $i$.
If $\chi=1$ and $s_t=0$, the positive donor row $t$ satisfies
$Z=D_t>d_0=.52$.  Since $k_s-\kappa_s(q)\ge0$,
Equation~\eqref{eq:scalar-type-zero} gives
$I_{s_i}(0;q,Z)\le-H_0$, which covers its incoming charge.
In the other cases the charge is either zero or $H_1=.25$.
Equation~\eqref{eq:scalar-type-one}, together with $q\le1$,
gives
\[
 I_{s_i}(0;q,Z)\le b-1.6z-C+k_s\le-.25,
\]
which covers these cases as well.  Convexity proves the pointwise
bound throughout $[0,c_*]$.  Apply it at $D=D_i$ and average over
$t$ in Equation~\eqref{eq:centered-row-bound} to obtain
Equation~\eqref{eq:incoming-charge}.

Finally, the incoming charge is $\Gamma_{ti}$ by definition, and
exchanging the two finite sums requires no symmetry of $q$.
Equation~\eqref{eq:outgoing-charge} bounds all
positive rows by their outgoing charges;
Equation~\eqref{eq:incoming-charge} makes all receiving rows pay
their incoming charges; all other rows have nonpositive contribution.
Every nonzero $\Gamma_{it}$ has $i\in\mathcal P$ and
$t\in\mathcal S$.  Exchanging the two indices in the finite sum with
weights $\pi_i\pi_t$ therefore gives
\[
 \E_iF_i\le\E_{i,t}\Gamma_{it}-\E_{i,t}\Gamma_{ti}=0.
\]
This proves the stronger bound $2z$ after the parameter reduction and,
since $z\le\Theta$, proves Equation~\eqref{eq:finite-inequality}.
\end{proof}
